
\chapter{ПОД/ФТ}\label{ch:aml}
ПОД/ФТ -- противодействие легализации (отмыванию) доходов
и финансированию терроризма; в~международной практике
для близкого набора задач используется аббревиатура AML
(anti-money laundering). Платёжный комплаенс по ПОД/ФТ
собирается вокруг 115-ФЗ
и~подзаконных актов Банка России и~Росфинмониторинга. Международные стандарты заданы FATF,
а~для СНГ её методологическую функцию исполняет ЕАГ.

\section{115-ФЗ и платёжный бизнес}\label{sec:aml-115fz}

115-ФЗ, Федеральный закон \enquote{О~противодействии
  легализации (отмыванию) доходов, полученных преступным путём, и
финансированию терроризма}, образует основу контроля ПОД/ФТ в~России~\cite{fz-115}. \highlight{Закон требует от оператора знать клиента,
наблюдать операции и~при формальных основаниях блокировать средства}.
Пока клиент не прошёл проверку, его операцию нельзя считать
доверенной.

Первая обязанность -- идентификация, при которой закон различает
полную и упрощённую процедуру, а подзаконные акты Банка России
задают состав сведений и документов по клиенту, представителю,
выгодоприобретателю и бенефициарному владельцу~\cite{cbr-499p}.
Регистрационных данных не всегда достаточно: где-то хватает минимального
набора, а где-то нужны документы, внешние реестры и дополнительные
вопросы. Объём проверки зависит от типа клиента, операции и~уровня
риска, поэтому привязываться к~быстро меняющимся порогам и~частным
исключениям в~печатной версии бессмысленно. Подробно идентификация
физлиц и ТСП разобрана в~главе~\ref{ch:kyc-kyb}.

\highlight{Вторая обязанность -- мониторинг, при котором оператор замечает
операции, подпадающие под критерии обязательного контроля,
и отправляет сведения в~Росфинмониторинг}. Этим задача не
ограничивается: надо ловить и~те операции, которые формально
ещё не описаны в~законе, но выбиваются из нормального профиля
клиента по поведению или по контексту. Для этого нужны регламент,
поток данных, маршрут проверки и~понятное решение по результату,
поскольку формально настроенный контроль либо ничего не увидит,
либо захлебнётся в~ложных срабатываниях.

115-ФЗ разделяет два режима контроля операций. Обязательный
контроль по статье~6 применяется к~закрытому перечню операций.
Сюда входят операции по сумме, начиная с~порога 1\,000\,000~руб.\
для денежных операций с~наличными, и~по виду, когда операция
совершается лицами из перечня экстремистов либо когда речь идёт
о~недвижимости свыше порога. Факультативный контроль по статье~7
обязывает оператора выявлять необычные и подозрительные операции
на основании собственных правил внутреннего контроля. Сведения об операциях, подлежащих
обязательному контролю, и~о~подозрительных операциях направляются
в~Росфинмониторинг через личный кабинет в~виде формализованного
электронного сообщения, ФЭС, в~установленные 115-ФЗ и нормативными
актами Банка России и Росфинмониторинга
сроки~\cite{fz-115,cbr-860p,rosfinmonitoring-oes}.\footnote{Состав
показателей и форматы ФЭС регулируются приказом Росфинмониторинга
№ 18 от 08.02.2022 и связанной с~ним нормативной базой;
требования к~правилам внутреннего контроля кредитных организаций
с 13.08.2025 устанавливает Положение Банка России 860-П, которое
заменило 375-П. Точные реквизиты сверяются перед публикацией.}

Третья обязанность -- замораживание денежных средств и~иного
имущества. Если клиент или бенефициар попадает в~перечни лиц,
связанных с~терроризмом или экстремизмом, организация обязана
применить меры по замораживанию и исполнить предусмотренный законом
порядок информирования~\cite{fz-115}. Технология сопоставления имён и~эскалации
разобрана в~главе~\ref{ch:sanctions}.

Банк России при надзоре за оператором платёжной инфраструктуры
проверяет не только формальные правила, но и~наличие рабочей
системы ПОД/ФТ, обученных сотрудников и~понятного маршрута
обработки случаев. Этот контур
определяет, кого продукт впускает
и~на каких условиях продолжает работу.

\section{Разбор QIWI: уроки для операционного риска}\label{sec:aml-qiwi}

21 февраля 2024 года Банк России отозвал лицензию
у~КИВИ Банка приказом № ОД-266~\cite{cbr-qiwi-license-2024}. В~официальном пресс-релизе
регулятор указал на систематические нарушения требований
законодательства по ПОД/ФТ, вовлечённость банка в~высокорисковые
операции, включая переводы в~пользу криптообменников и нелегальных
онлайн-казино, а~также на случаи открытия кошельков QIWI с
использованием персональных данных граждан без их ведома.

В~инженерных терминах нарушения выглядят так.
115-ФЗ требует от банков идентифицировать клиентов и проводить
контроль операций; открытие кошельков без ведома граждан означает,
что процедура идентификации либо не выполнялась, либо не
проверялась системой контроля. Переводы в~пользу криптообменников
и~казино становятся проблемой тогда, когда transaction monitoring
не маркирует их как высокорисковые или маркирует, но событие
не попадает на ручной разбор. Оба сбоя -- технические: отсутствие
автоматического триггера для enhanced due diligence и неполный
audit trail по спорным транзакциям.

Компания, зависящая от одного провайдера, наследует все его
слабые места и~в~момент отзыва лицензии теряет рабочий платёжный
канал.

Рост объёма без роста контроля создаёт слепые пятна. Регуляторные
предписания стоит читать как предупреждение.

\importantnote{%
  Для бизнеса с~зависимостью от одного платёжного
  провайдера нужен заранее проверенный план непрерывности:
  альтернативный провайдер, резервные счета и~маршрут переключения.

}
\section{FATF и~ЕАГ для СНГ}\label{sec:aml-fatf-eag}

FATF задаёт международные стандарты через рекомендации, взаимные
оценки и публичные списки юрисдикций под усиленным
наблюдением~\cite{fatf-recs,fatf-increased-monitoring}. С 24~февраля
2023~года членство Российской Федерации в~FATF
приостановлено~\cite{fatf-russia-suspension-2023}: формально это
означает, что Россия не участвует в~выработке решений FATF,
не направляет своих представителей в~её рабочие группы и~не
подвергается очередным взаимным оценкам в~стандартном порядке.
В~СНГ функцию региональной FATF-style body исполняет Евразийская
группа по противодействию легализации преступных доходов
и финансированию терроризма, ЕАГ. Именно ЕАГ
проводит взаимные оценки государств региона и публикует
методологические документы, на которые опираются национальные
регуляторы стран
СНГ~\cite{eag-overview,eag-mutual-evaluations}.\footnote{Точная
дата приостановки членства РФ в~FATF и~текущий статус
взаимодействия, приостановлено или исключено, сверяются по сайту
FATF и официальным сообщениям Росфинмониторинга перед публикацией.}

После приостановки членства стандарты FATF продолжают влиять на
платёжный бизнес в~России косвенно. Банки-корреспонденты и платёжные
посредники в~третьих юрисдикциях смотрят на риск через рекомендации
FATF и публичные списки. Крупные банки практикуют политику снижения
контрагентского риска, de-risking, и сокращают или полностью
прекращают отношения с~целыми классами клиентов или
юрисдикциями~\cite{fatf-correspondent-banking}. Это означает прежде всего трудности открытия и поддержания
корреспондентских счетов в~валютах недружественных стран,
поскольку появляются дополнительные проверки и ограничивается
доступность канала.

\practicenote{%
  Перед выходом в~новую юрисдикцию проверяются оценка ЕАГ
  для стран СНГ либо FATF для остальных, последняя публичная
  взаимная оценка страны и~местная практика ПОД/ФТ. Если что-то
  из этого вызывает сомнение, заранее закладывается запас
  по времени и~каналам расчёта.

}
